\documentclass[11pt]{article}

\usepackage[a4paper,margin=1in]{geometry}
\usepackage{amsmath,amssymb,amsthm,mathtools}
\usepackage{microtype}
\usepackage[hidelinks]{hyperref}

\newtheorem{theorem}{Theorem}[section]
\newtheorem{proposition}[theorem]{Proposition}
\newtheorem{lemma}[theorem]{Lemma}
\newtheorem{corollary}[theorem]{Corollary}
\theoremstyle{definition}
\newtheorem{definition}[theorem]{Definition}
\theoremstyle{remark}
\newtheorem{remark}[theorem]{Remark}

\newcommand{\Hom}{\operatorname{Hom}}
\newcommand{\homn}{\operatorname{hom}}
\newcommand{\tr}{\operatorname{Tr}}
\newcommand{\R}{\mathbb{R}}
\newcommand{\one}{\mathbf{1}}
\newcommand{\e}{\mathrm{e}}

\title{Even-Anchored Interpolation for Odd Cycles\\
and Domination Exponents}
\author{Taeyoung Kim}
\date{\today}
\hypersetup{
  pdftitle={Even-Anchored Interpolation for Odd Cycles and Domination Exponents},
  pdfauthor={Taeyoung Kim},
  pdfsubject={Cycle homomorphism inequalities and odd-cycle domination exponents}
}

\begin{document}
\maketitle

\begin{abstract}
Let $T$ be a finite simple graph.  We prove that, whenever $b$ is even,
$n,m$ are odd, and $2\le b<n<m$,
\[
  t(C_b,T)^{m-n}t(C_m,T)^{n-b}
  \ge t(C_n,T)^{m-b},
\]
where $C_2$ denotes $K_2$.  The proof roots the
three closed-walk counts at a vertex, applies Sa\u{g}lam's same-parity
power inequality to $A^{b/2}e_x$, and sums the resulting pointwise
inequality by H\"older's inequality.  This answers Problem~6 of
Blekherman--Raymond--Razborov--Wei.  Combining the edge-anchored case
with their pruning argument also gives, for all $k>\ell\ge1$,
\[
  C(C_{2k+1},C_{2\ell+1})
  \le
  \frac{2k(2k-2\ell+1)-1}
       {2\ell(2k-2\ell+1)-1},
\]
extending their second upper-bound formula to the full parameter range.
\end{abstract}

\section{Definitions and main results}

All graphs are finite, simple, undirected, and loopless.  For finite
graphs $F,T$, write
\[
  \homn(F,T):=|\Hom(F,T)|,
  \qquad
  t(F,T):=\frac{\homn(F,T)}{|V(T)|^{|V(F)|}}.
\]
For a uniform treatment of the edge case, set
\[
  C_2:=K_2,
\]
while $C_r$ has its usual meaning for $r\ge3$.  Thus $C_2$ is a single
edge, not a two-edge multigraph.

Our main interpolation theorem is the following.

\begin{theorem}[Even-anchored interpolation for odd cycles]
  \label{thm:interpolation}
  Let $b,n,m$ be integers such that
  \[
    2\le b<n<m,
    \qquad
    b\equiv0\pmod2,
    \qquad
    n\equiv m\equiv1\pmod2.
  \]
  Then every finite simple graph $T$ satisfies
  \begin{equation}
    \homn(C_b,T)^{m-n}\homn(C_m,T)^{n-b}
    \ge
    \homn(C_n,T)^{m-b}.
    \label{eq:main-hom}
  \end{equation}
  Equivalently,
  \begin{equation}
    t(C_b,T)^{m-n}t(C_m,T)^{n-b}
    \ge
    t(C_n,T)^{m-b}.
    \label{eq:main-density}
  \end{equation}
\end{theorem}

Taking consecutive odd cycles gives precisely the inequality posed in
\cite[Problem~6]{BRRW2025}.

\begin{corollary}[Problem~6 of \cite{BRRW2025}]
  \label{cor:problem6}
  Let $i>j\ge1$.  Then every finite simple graph $T$ satisfies
  \begin{equation}
    t(C_{2j},T)^2
    t(C_{2i+1},T)^{2i-1-2j}
    \ge
    t(C_{2i-1},T)^{2i+1-2j},
    \label{eq:problem6}
  \end{equation}
  where $C_2$ denotes $K_2$.
\end{corollary}

We also obtain an improved upper bound for odd-cycle domination
exponents.

\begin{definition}[Homomorphism-density domination exponent]
  For graphs $G,H$, let
  \[
    C(G,H)
    :=
    \min\bigl\{c\ge0:
    t(G,T)\ge t(H,T)^c
    \text{ for every finite simple graph }T\bigr\},
  \]
  whenever this set is nonempty.
\end{definition}

\begin{theorem}[Odd-cycle domination bound]
  \label{thm:domination}
  For all integers $k>\ell\ge1$,
  \begin{equation}
    C(C_{2k+1},C_{2\ell+1})
    \le
    \frac{2k(2k-2\ell+1)-1}
         {2\ell(2k-2\ell+1)-1}.
    \label{eq:domination}
  \end{equation}
\end{theorem}

\begin{corollary}[Adjacent odd cycles]
  \label{cor:adjacent}
  For every $i\ge2$,
  \[
    C(C_{2i+1},C_{2i-1})
    \le \frac{6i-1}{6i-7}.
  \]
\end{corollary}

The proof of Theorem~\ref{thm:interpolation} uses one non-elementary
input, Sa\u{g}lam's same-parity power inequality.  The proof of
Theorem~\ref{thm:domination} then combines the case $b=2$ with an
explicit version of the pruning and tensor-power argument from
\cite{BRRW2025}.

\section{Rooted interpolation}

We first record the matrix inequality that drives the proof.

\begin{theorem}[Sa\u{g}lam \cite{Saglam2018}]
  \label{thm:saglam}
  Let $\Omega$ be a finite set, let
  $S\in\R^{\Omega\times\Omega}$ be symmetric and entrywise
  nonnegative, and let $u,v\in\R^\Omega$ be entrywise nonnegative unit
  vectors.  If $K\ge R$ are positive integers of the same parity, then
  \begin{equation}
    \langle v,S^K u\rangle^R
    \ge
    \langle v,S^R u\rangle^K.
    \label{eq:saglam}
  \end{equation}
\end{theorem}

Let $T$ have $N\ge1$ vertices, and let $A$ be its adjacency matrix.
For $x\in V(T)$ and $r\ge0$, put
\[
  c_r(x):=(A^r)_{xx}=\langle e_x,A^r e_x\rangle.
\]
For $r\ge1$, this is the number of closed walks of length $r$ based at
$x$.

\begin{lemma}[Trace and rooted-walk identities]
  \label{lem:rooted-identities}
  For every $r\ge2$,
  \begin{equation}
    \homn(C_r,T)=\tr(A^r)=\sum_{x\in V(T)}c_r(x).
    \label{eq:trace-hom}
  \end{equation}
  Moreover, if $b\ge2$ is even and
  \[
    w_x:=A^{b/2}e_x,
  \]
  then $w_x$ is entrywise nonnegative and, for every $r\ge b$,
  \begin{equation}
    \|w_x\|_2^2=c_b(x),
    \qquad
    \langle w_x,A^{r-b}w_x\rangle=c_r(x).
    \label{eq:anchor-identities}
  \end{equation}
\end{lemma}

\begin{proof}
  For $r\ge3$, expansion of the diagonal entries of $A^r$ identifies
  $\tr(A^r)$ with the number of closed walks of length $r$, hence with
  $\homn(C_r,T)$.  For $r=2$, the fact that $A$ is a symmetric
  $0$--$1$ matrix gives
  \[
    \tr(A^2)
    =\sum_{x,y}A_{xy}A_{yx}
    =\sum_{x,y}A_{xy}
    =2|E(T)|
    =\homn(K_2,T).
  \]
  This proves \eqref{eq:trace-hom}.

  Since every power of the entrywise nonnegative matrix $A$ is
  entrywise nonnegative, so is $w_x$.  Symmetry of $A$ gives
  \[
    \|w_x\|_2^2
    =\langle e_x,A^b e_x\rangle=c_b(x),
  \]
  and, for $r\ge b$,
  \[
    \langle w_x,A^{r-b}w_x\rangle
    =\langle e_x,A^{b/2}A^{r-b}A^{b/2}e_x\rangle
    =\langle e_x,A^r e_x\rangle=c_r(x).
  \]
\end{proof}

\begin{lemma}[Rooted even-anchor interpolation]
  \label{lem:rooted-interpolation}
  Under the hypotheses of Theorem~\ref{thm:interpolation}, for every
  $x\in V(T)$,
  \begin{equation}
    c_b(x)^{m-n}c_m(x)^{n-b}
    \ge c_n(x)^{m-b}.
    \label{eq:rooted-interpolation}
  \end{equation}
\end{lemma}

\begin{proof}
  Fix $x$ and let $w_x=A^{b/2}e_x$.  If $c_b(x)=0$, then
  \eqref{eq:anchor-identities} gives $\|w_x\|_2=0$, so $w_x=0$.
  Consequently $c_n(x)=c_m(x)=0$, and
  \eqref{eq:rooted-interpolation} is immediate.

  Suppose now that $c_b(x)>0$ and set
  \[
    u_x:=\frac{w_x}{\|w_x\|_2},
    \qquad
    K:=m-b,
    \qquad
    R:=n-b.
  \]
  The vector $u_x$ is an entrywise nonnegative unit vector.  Since $b$
  is even and $n,m$ are odd, $K>R$ are positive odd integers.
  Theorem~\ref{thm:saglam}, applied with $S=A$ and $u=v=u_x$, yields
  \[
    \langle u_x,A^{m-b}u_x\rangle^{n-b}
    \ge
    \langle u_x,A^{n-b}u_x\rangle^{m-b}.
  \]
  By \eqref{eq:anchor-identities},
  \[
    \langle u_x,A^{r-b}u_x\rangle
    =\frac{c_r(x)}{c_b(x)}
    \qquad(r\in\{n,m\}).
  \]
  Substitution and multiplication by the positive number
  $c_b(x)^{m-b}$ give
  \[
    c_b(x)^{(m-b)-(n-b)}c_m(x)^{n-b}
    \ge c_n(x)^{m-b},
  \]
  which is \eqref{eq:rooted-interpolation}.
\end{proof}

\begin{proof}[Proof of Theorem~\ref{thm:interpolation}]
  Define
  \[
    \alpha:=\frac{m-n}{m-b},
    \qquad
    \beta:=\frac{n-b}{m-b}.
  \]
  Then $\alpha,\beta\in(0,1)$ and $\alpha+\beta=1$.
  Taking the $(m-b)$-th root in
  \eqref{eq:rooted-interpolation} gives
  \[
    c_n(x)\le c_b(x)^\alpha c_m(x)^\beta.
  \]
  H\"older's inequality with conjugate exponents $1/\alpha$ and
  $1/\beta$ therefore implies
  \[
  \begin{aligned}
    \sum_xc_n(x)
    &\le \sum_x c_b(x)^\alpha c_m(x)^\beta\\
    &\le
    \left(\sum_xc_b(x)\right)^\alpha
    \left(\sum_xc_m(x)\right)^\beta.
  \end{aligned}
  \]
  Using \eqref{eq:trace-hom} and raising to the power $m-b$ proves
  \eqref{eq:main-hom}.

  To pass to densities, substitute
  $\homn(F,T)=N^{|V(F)|}t(F,T)$.  The powers of $N$ cancel because
  \[
    b(m-n)+m(n-b)=n(m-b).
  \]
  This proves \eqref{eq:main-density}.
\end{proof}

\begin{proof}[Proof of Corollary~\ref{cor:problem6}]
  Apply Theorem~\ref{thm:interpolation} with
  \[
    b=2j,
    \qquad
    n=2i-1,
    \qquad
    m=2i+1.
  \]
  The assumption $i>j$ gives $2j<2i-1$, and the three exponents in
  \eqref{eq:main-density} become
  \[
    m-n=2,
    \qquad
    n-b=2i-1-2j,
    \qquad
    m-b=2i+1-2j.
  \]
  This is exactly \eqref{eq:problem6}.
\end{proof}

\begin{remark}[How H\"older removes regularity]
  For $b=2$, the rooted inequality reads
  \[
    \deg_T(x)^{m-n}c_m(x)^{n-2}
    \ge c_n(x)^{m-2}.
  \]
  In a regular graph the degree is constant and can be pulled out of
  the sum.  For an arbitrary graph, H\"older's inequality performs the
  same summation exactly:
  \[
    \sum_x\deg_T(x)^\alpha c_m(x)^\beta
    \le
    \left(\sum_x\deg_T(x)\right)^\alpha
    \left(\sum_xc_m(x)\right)^\beta.
  \]
\end{remark}

\section{Domination exponents}

We now prove Theorem~\ref{thm:domination}.  The argument follows the
pruning framework of \cite[Theorem~4.6]{BRRW2025}, with constants and
normalizations included for completeness.

\subsection{Even cycles and rooted edges}

\begin{lemma}[Even-cycle lower bound]
  \label{lem:even-cycle}
  For every $r\ge2$ and every finite simple graph $T$,
  \begin{equation}
    t(C_{2r},T)\ge t(K_2,T)^{2r}.
    \label{eq:even-cycle}
  \end{equation}
\end{lemma}

\begin{proof}
  Let $A$ be the adjacency matrix of $T$, let $N=|V(T)|$, and put
  $p=t(K_2,T)$.  If $\lambda_{\max}$ is the largest eigenvalue of $A$,
  the Rayleigh--Ritz principle applied to $\one$ gives
  \[
    \lambda_{\max}
    \ge
    \frac{\langle\one,A\one\rangle}{\langle\one,\one\rangle}
    =Np.
  \]
  If $\lambda_1,\ldots,\lambda_N$ are the eigenvalues of $A$, then
  \[
    \homn(C_{2r},T)
    =\tr(A^{2r})
    =\sum_{s=1}^N\lambda_s^{2r}
    \ge\lambda_{\max}^{2r}
    \ge(Np)^{2r}.
  \]
  Division by $N^{2r}$ proves \eqref{eq:even-cycle}.
\end{proof}

Give $C_q$ the cyclic vertex set $\mathbb{Z}/q\mathbb{Z}$ and
distinguish the oriented edge $(0,1)$.  For an oriented edge $(u,v)$
of $T$, define
\[
  h_q^T(u,v)
  :=
  \bigl|\{\phi\in\Hom(C_q,T):
  \phi(0)=u,\ \phi(1)=v\}\bigr|.
\]
Write
\[
  \vec E(T)
  :=\{(u,v)\in V(T)^2:\{u,v\}\in E(T)\}.
\]
Partitioning homomorphisms according to the image of $(0,1)$ gives
\begin{equation}
  \sum_{(u,v)\in\vec E(T)}h_q^T(u,v)=\homn(C_q,T).
  \label{eq:rooted-edge-sum}
\end{equation}
Reflection of $C_q$ across the distinguished edge also gives
\begin{equation}
  h_q^T(u,v)=h_q^T(v,u).
  \label{eq:rooted-edge-symmetry}
\end{equation}

\begin{lemma}[Pruning to cycle-rich edges]
  \label{lem:pruning}
  Fix $q\ge3$.  Let $T_0$ be a graph on $N\ge3$ vertices and suppose
  \[
    H_0:=\homn(C_q,T_0)>0.
  \]
  Then $T_0$ has a spanning subgraph $T$ such that, with
  \[
    H:=\homn(C_q,T),
    \qquad
    M:=|E(T)|,
  \]
  one has
  \begin{equation}
    H\ge\e^{-3/2}H_0
    \label{eq:pruning-preservation}
  \end{equation}
  and, for every $(u,v)\in\vec E(T)$,
  \begin{equation}
    h_q^T(u,v)>
    \frac{H}{8qM\log N}.
    \label{eq:uniform-rooted-edge}
  \end{equation}
  Here and below $\log$ denotes the natural logarithm.
\end{lemma}

\begin{proof}
  Starting from $T_0$, suppose the current graph $T_i$ has $i$ edges,
  and write $H_i=\homn(C_q,T_i)$.  For $e\in E(T_i)$, let $L_i(e)$ be
  the number of homomorphisms $C_q\to T_i$ whose image uses $e$.
  If some edge satisfies
  \[
    L_i(e)\le\frac{H_i}{4i\log N},
  \]
  delete it.  The resulting graph $T_{i-1}$ satisfies
  \begin{equation}
    H_{i-1}=H_i-L_i(e)
    \ge H_i\left(1-\frac1{4i\log N}\right).
    \label{eq:pruning-step}
  \end{equation}
  Stop when no such edge remains, and call the terminal graph $T$.
  Since every factor in \eqref{eq:pruning-step} is positive, all
  intermediate $H_i$ are positive; in particular, $T$ has at least one
  edge.

  Let $M_0=|E(T_0)|$.  Iterating \eqref{eq:pruning-step} gives
  \[
    H\ge
    H_0\prod_{i=M+1}^{M_0}
    \left(1-\frac1{4i\log N}\right).
  \]
  Since $N\ge3$, each $x_i=1/(4i\log N)$ lies in $[0,1/2]$.
  Using $\log(1-x)\ge-2x$ on this interval,
  \[
  \begin{aligned}
    \log\prod_{i=M+1}^{M_0}(1-x_i)
    &\ge-\frac1{2\log N}\sum_{i=1}^{M_0}\frac1i\\
    &\ge-\frac{1+\log M_0}{2\log N}
    \ge-\frac32.
  \end{aligned}
  \]
  The last inequality uses $M_0<N^2$ and $\log N\ge1$.  This proves
  \eqref{eq:pruning-preservation}.

  At termination, every edge $e=\{u,v\}$ satisfies
  \[
    L_T(e)>\frac{H}{4M\log N}.
  \]
  For $a\in\mathbb{Z}/q\mathbb{Z}$, let $I_a(e)$ count those
  homomorphisms for which the source edge $\{a,a+1\}$ maps to $e$.
  Every homomorphism counted by $L_T(e)$ contributes to at least one
  $I_a(e)$, so
  \[
    \sum_a I_a(e)\ge L_T(e).
  \]
  Rotation symmetry makes all $I_a(e)$ equal, while
  \eqref{eq:rooted-edge-symmetry} gives
  \[
    I_0(e)=h_q^T(u,v)+h_q^T(v,u)=2h_q^T(u,v).
  \]
  Hence
  \[
    2q\,h_q^T(u,v)
    =\sum_aI_a(e)
    >\frac{H}{4M\log N},
  \]
  which proves \eqref{eq:uniform-rooted-edge}.
\end{proof}

\begin{lemma}[Gluing two edge-rooted cycles]
  \label{lem:gluing}
  Let $q,h\ge3$ and set $m=q+h-2$.  Suppose that a graph $T$ on
  $N\ge3$ vertices satisfies \eqref{eq:uniform-rooted-edge} for $C_q$.
  Put
  \[
    a:=t(C_q,T),
    \qquad
    z:=t(C_h,T),
    \qquad
    p:=t(K_2,T),
    \qquad
    y:=t(C_m,T).
  \]
  Then
  \begin{equation}
    y>\frac1{4q\log N}\frac{az}{p}.
    \label{eq:gluing}
  \end{equation}
\end{lemma}

\begin{proof}
  The hypothesis implies $E(T)\ne\varnothing$, so $p>0$.  Label the
  vertices of $C_m$ cyclically by $0,\ldots,m-1$, and add the chord
  $\{0,q-1\}$.  The resulting graph $C_m^+$ is the union of a $q$-cycle
  and an $h$-cycle along that chord.  Consequently,
  \[
  \begin{aligned}
    \homn(C_m,T)
    &\ge\homn(C_m^+,T)\\
    &=\sum_{(u,v)\in\vec E(T)}
      h_q^T(u,v)h_h^T(u,v)\\
    &>
    \frac{\homn(C_q,T)}
         {8q|E(T)|\log N}
    \sum_{(u,v)\in\vec E(T)}h_h^T(u,v)\\
    &=
    \frac{\homn(C_q,T)\homn(C_h,T)}
         {8q|E(T)|\log N}.
  \end{aligned}
  \]
  Now $|E(T)|=N^2p/2$ and $q+h=m+2$.  Dividing by $N^m$ therefore
  gives \eqref{eq:gluing}.
\end{proof}

\subsection{Combining the inequalities}

Fix $k>\ell\ge1$ and put
\[
  \Delta:=k-\ell,
  \qquad
  n:=2\ell+1,
  \qquad
  m:=2k+1,
  \qquad
  h:=2\Delta+2.
\]
Thus $n+h-2=m$.

\begin{lemma}[Polylogarithmic domination after pruning]
  \label{lem:polylog}
  There are constants $c_0>0$ and $\gamma\ge0$, depending only on
  $k,\ell$, such that every graph $T_0$ on $N\ge3$ vertices satisfies
  \begin{equation}
    t(C_m,T_0)
    \ge
    c_0(\log N)^{-\gamma}
    t(C_n,T_0)^\eta,
    \label{eq:polylog}
  \end{equation}
  where
  \begin{equation}
    \eta
    :=
    \frac{2k(2\Delta+1)-1}
         {2\ell(2\Delta+1)-1}.
    \label{eq:eta}
  \end{equation}
\end{lemma}

\begin{proof}
  Write
  \[
    a_0:=t(C_n,T_0),
    \qquad
    y_0:=t(C_m,T_0).
  \]
  If $a_0=0$, there is nothing to prove.  Suppose $a_0>0$.
  Lemma~\ref{lem:pruning}, applied with $q=n$, gives a spanning
  subgraph $T\subseteq T_0$ satisfying
  \eqref{eq:pruning-preservation} and
  \eqref{eq:uniform-rooted-edge}.  Set
  \[
    a:=t(C_n,T),
    \qquad
    y:=t(C_m,T),
    \qquad
    p:=t(K_2,T).
  \]
  Since the subgraph is spanning,
  \begin{equation}
    a\ge\e^{-3/2}a_0,
    \qquad
    y_0\ge y.
    \label{eq:preserved-densities}
  \end{equation}

  Apply Lemma~\ref{lem:gluing} with $q=n$ and
  $h=2\Delta+2$, and then apply Lemma~\ref{lem:even-cycle} to $C_h$.
  This gives
  \begin{equation}
    y>
    \frac1{4n\log N}\frac{a\,t(C_h,T)}p
    \ge
    \frac1{4n\log N}a p^{2\Delta+1}.
    \label{eq:gluing-lower}
  \end{equation}
  On the other hand, Theorem~\ref{thm:interpolation} with
  $(b,n,m)=(2,2\ell+1,2k+1)$ gives the exact constraint
  \begin{equation}
    p^{2\Delta}y^{2\ell-1}\ge a^{2k-1}.
    \label{eq:anchor-constraint}
  \end{equation}

  Introduce
  \[
    A_0:=2\ell-1,
    \quad
    B_0:=2\Delta,
    \quad
    C_0:=2k-1=A_0+B_0,
    \quad
    D_0:=2\Delta+1=B_0+1.
  \]
  Since $a>0$, \eqref{eq:anchor-constraint} implies $p,y>0$.
  Raising \eqref{eq:gluing-lower} to $B_0$ and
  \eqref{eq:anchor-constraint} to $D_0$ gives
  \[
    y^{B_0}>
    (4n\log N)^{-B_0}a^{B_0}p^{B_0D_0},
    \qquad
    p^{B_0D_0}
    \ge\frac{a^{C_0D_0}}{y^{A_0D_0}}.
  \]
  Eliminating $p$ yields
  \[
    y^{B_0+A_0D_0}>
    (4n\log N)^{-B_0}a^{B_0+C_0D_0}.
  \]
  Therefore
  \begin{equation}
    y>
    (4n)^{-\theta}(\log N)^{-\theta}a^\eta,
    \qquad
    \theta:=\frac{B_0}{B_0+A_0D_0},
    \label{eq:combined}
  \end{equation}
  where
  \[
  \begin{aligned}
    B_0+A_0D_0
    &=(A_0+1)(B_0+1)-1
      =2\ell(2\Delta+1)-1,\\
    B_0+C_0D_0
    &=(A_0+B_0+1)(B_0+1)-1
      =2k(2\Delta+1)-1.
  \end{aligned}
  \]
  Thus the exponent in \eqref{eq:combined} is exactly the number
  $\eta$ in \eqref{eq:eta}.  Finally,
  \eqref{eq:preserved-densities} and \eqref{eq:combined} give
  \[
    y_0>
    (4n)^{-\theta}\e^{-3\eta/2}
    (\log N)^{-\theta}a_0^\eta.
  \]
  Hence \eqref{eq:polylog} holds with
  \[
    c_0=(4n)^{-\theta}\e^{-3\eta/2},
    \qquad
    \gamma=\theta.
  \]
\end{proof}

\begin{lemma}[Tensor-power removal]
  \label{lem:tensor}
  Let $H_1,H_2$ be fixed graphs and let $c>0$.  Suppose that constants
  $K>0$ and $D\ge0$ satisfy
  \[
    t(H_1,T)\ge
    K(\log |V(T)|)^{-D}t(H_2,T)^c
  \]
  for every graph $T$ with at least three vertices.  Then
  \[
    t(H_1,T)\ge t(H_2,T)^c
  \]
  for every such graph $T$.
\end{lemma}

\begin{proof}
  Fix $T$ with $N\ge3$ vertices and let $T^{\times r}$ be its $r$-fold
  categorical tensor product.  Coordinatewise factorization of
  homomorphisms gives
  \[
    t(H,T^{\times r})=t(H,T)^r
  \]
  for every fixed graph $H$.  Applying the assumed inequality to
  $T^{\times r}$ gives
  \[
    t(H_1,T)^r
    \ge
    K(r\log N)^{-D}t(H_2,T)^{cr}.
  \]
  Taking $r$-th roots and sending $r\to\infty$, we use
  \[
    K^{1/r}\longrightarrow1,
    \qquad
    (r\log N)^{-D/r}\longrightarrow1
  \]
  to obtain the result.
\end{proof}

\begin{proof}[Proof of Theorem~\ref{thm:domination}]
  Lemma~\ref{lem:polylog} and Lemma~\ref{lem:tensor} give
  \[
    t(C_{2k+1},T)\ge t(C_{2\ell+1},T)^\eta
  \]
  for every target $T$ with at least three vertices, where
  \[
    \eta
    =
    \frac{2k(2k-2\ell+1)-1}
         {2\ell(2k-2\ell+1)-1}.
  \]
  A simple graph with at most two vertices is bipartite, so it admits
  no homomorphism from an odd cycle.  Both sides are then zero.
  Thus the displayed inequality holds for every finite simple target,
  and the definition of $C$ proves \eqref{eq:domination}.
\end{proof}

\begin{proof}[Proof of Corollary~\ref{cor:adjacent}]
  Set $k=i$ and $\ell=i-1$ in Theorem~\ref{thm:domination}.  Then
  $2k-2\ell+1=3$, so the numerator is $6i-1$ and the denominator is
  $6i-7$.
\end{proof}

\section{Remarks}

Blekherman--Raymond--Razborov--Wei prove
\[
  C(C_{2k+1},C_{2\ell+1})
  \le
  \begin{cases}
    \dfrac{2(k+1)}{2\ell+1},&k\le2\ell-1,\\[2mm]
    \dfrac{2k(2k-2\ell+1)-1}
          {2\ell(2k-2\ell+1)-1},&k\ge2\ell.
  \end{cases}
\]
Theorem~\ref{thm:domination} extends the second expression to every
$k>\ell$.  If $D=2k-2\ell+1$, then
\[
  \frac{2(k+1)}{2\ell+1}
  -
  \frac{2kD-1}{2\ell D-1}
  =
  \frac{(4\ell-2k-1)D}
       {(2\ell+1)(2\ell D-1)}.
\]
Thus the new bound is strictly stronger throughout the former
first-branch range $k\le2\ell-1$.  Their projective-plane construction
gives the lower bound
\[
  C(C_{2k+1},C_{2\ell+1})
  \ge\frac{4k^2-1}{4k\ell-1},
\]
and the remaining gap is
\[
  \frac{2(k-\ell)(2\ell-1)}
  {(4k\ell-1)\bigl(2\ell(2k-2\ell+1)-1\bigr)}.
\]

The $b=2$ case of Theorem~\ref{thm:interpolation} is specifically a
finite simple-target statement.  Indeed, for a $0$--$1$ adjacency
matrix,
\[
  \tr(A^2)=\homn(K_2,T),
\]
whereas for a graphon $W$,
\[
  \tr(T_W^2)
  =\int_{[0,1]^2}W(x,y)^2\,dx\,dy
\]
need not equal the edge density $\int W$.  A direct operator analogue
therefore has the two-edge multigraph density, rather than
$t(K_2,W)$, as its length-two anchor.

\begin{thebibliography}{9}

\bibitem{Saglam2018}
M.~Sa\u{g}lam,
\emph{Near log-convexity of measured heat in (discrete) time and consequences},
arXiv:1808.06717, 2018.

\bibitem{BRRW2025}
G.~Blekherman, A.~Raymond, A.~Razborov, and F.~Wei,
\emph{On domination exponents for pairs of graphs},
arXiv:2506.12151v2, 2025.

\end{thebibliography}

\end{document}
